\documentclass[11pt,a4paper]{article}

\usepackage[T1]{fontenc}
\usepackage[utf8]{inputenc}
\usepackage[english]{babel}
\usepackage[a4paper,margin=2.8cm]{geometry}
\usepackage{microtype}
\usepackage[hidelinks]{hyperref}

\hypersetup{
  pdftitle={The Aleph},
  pdfauthor={Juan José Gómez Cadenas}
}

\setlength{\parindent}{1.25em}
\setlength{\parskip}{0.35em}
\raggedbottom

\title{The Aleph}
\author{Juan José Gómez Cadenas}
\date{}

\begin{document}

\maketitle

It was strange to read \emph{The Aleph} again. I had not done it in decades, and I had the feeling one has when meeting a friend again, perhaps an old flame, whom one has not seen in years. I was excited, slightly apprehensive, full of hope. One never approaches Borges lightly; there is always the touch of the sacred floating around. And yet, I was afraid that the reality of the text would not live up to my exalted memories.

It was even better than I remembered, like being received with a warm embrace, an unexpected kiss that leaves on the lips the flavour of mistela.

\emph{The Aleph} is, primarily, a story of tragic love:

\begin{quotation}
On the burning February morning Beatriz Viterbo died, after braving an agony that never for a single moment gave way to self-pity or fear, I noticed that the sidewalk billboards around Constitution Plaza were advertising some new brand or other of American cigarettes. The fact pained me, for I realised that the wide and ceaseless universe was already slipping away from her and that this slight change was the first of an endless series.\footnote{Quotations from Borges are taken from the translation by Norman Thomas di Giovanni in collaboration with the author, \emph{The Aleph and Other Stories, 1933--1969} (New York: Dutton, 1970). Daneri's stanza, which that version omits, is my own translation.}
\end{quotation}

It was Borges himself who recommended beginning a good story with a birth, a death, or a stranger who arrives in a nameless town. And the death of Beatriz, in the first line, not only sets the sombre tone of the story but also refers the reader to that other Beatrice, whom Dante exalted in his \emph{Comedy}. But there is more. The next paragraph informs us of an imperceptible change that contributes to steering the universe away from the deceased. Only from her? Borges subtly invokes the reality of the constant change of the cosmos, which in turn reflects the incessant increase of entropy. Since the death of Beatriz Viterbo, in 1929, that stubbornly growing entropy has been changing the universe ---with a cigarette advertisement, with a world war, with man's arrival on the moon, with the recent irruption of AI---, but above all, it has been disordering it. Almost a century has passed since the events Borges narrates. In the Milky Way, some hundred stars will have gone out, more than a trillion if we count the whole universe. Since the death of Beatriz, nearly a thousand galaxies have emitted the last photon that will ever reach us. Borges intuits ---and captures, with a trivial detail--- a cosmic tragedy. Our universe is expanding, pushed by dark energy, accelerating ever faster, cooling little by little, towards an inescapable heat death.

There follows, after the opening paragraph, what we might call the burlesque part of the story, starring Beatriz's cousin, the insufferable Carlos Argentino Daneri, whom Borges describes with cruel irony:

\begin{quotation}
Beatriz had been tall, frail, slightly stooped; in her walk there was (if the oxymoron may be allowed) a kind of uncertain grace, a hint of expectancy. Carlos Argentino was pink-faced, overweight, gray-haired, fine-featured. […] His mental activity was continuous, deeply felt, far-ranging, and ---all in all--- meaningless.
\end{quotation}

From the death of Beatriz to the time in which the story takes place (1941), twelve years have passed, during which Borges regularly repeats his visits to the relative he cannot stand, in exchange for clinging to the little he has left of his beloved:

\begin{quotation}
Beatriz Viterbo in profile and in full colour; Beatriz wearing a mask, during the Carnival of 1921; Beatriz at her First Communion; Beatriz on the day of her wedding to Roberto Alessandri; Beatriz soon after her divorce, at a luncheon at the Turf Club; Beatriz at a seaside resort in Quilmes with Delia San Marco Porcel and Carlos Argentino; Beatriz with the Pekingese lapdog given her by Villegas Haedo; Beatriz, front and three-quarter views, smiling, hand on her chin…
\end{quotation}

The price he pays for those images is high. Carlos Argentino is the living embodiment of everything Borges hates, capable of boasting of stanzas such as this one:

\begin{verse}
On the right hand of the routine post \\
(coming, of course, from the North-northwest) \\
a skeleton is bored ---Colour? Celestewhite; \\
which gives the sheep pen the look of an ossuary.
\end{verse}

Four alexandrines that condense the folly of the undertaking the poetaster intends to carry out: a poem, entitled ``The Earth'', that aspires to describe the whole planet. Not content with his preposterous ambition, Argentino seems determined to exhaust every variation of bad taste:

\begin{quotation}
[…] where at first ``blue'' had been good enough, he now wallowed in ``azures'', ``ceruleans'', and ``ultramarines''. The word ``milky'' was too easy for him; in the course of an impassioned description of a shed where wool was washed, he chose such words as ``lacteal'', ``lactescent'', and even made one up ---``lactinacious''.
\end{quotation}

Borges (and the reader) tolerates Argentino with great difficulty, and the whole story seems headed towards a melancholy absurdity… Until suddenly, as if fallen from the sky, an Aleph appears.

Borges learns of this circumstance when Argentino informs him that his landlords plan to demolish the building where he lives, with the intention of enlarging a tearoom. For the writer, the news is a low blow, since with the destruction of the house the memory of Beatriz fades a little more (once again the heat death of the Universe). The fake poet, on the other hand, complains that, to finish his opus, the building is indispensable to him, since in a corner of the cellar an Aleph lies hidden. And he explains, in passing, that the Aleph is one of the points in space that contains all points, the place where all the places of the orb are, without blending, seen from every angle.

Here the story turns dreamlike. An incredulous Borges recounts the tantrum of the spoiled child, convinced that the Aleph in question, besides being his property, exists for the sole purpose of helping him compose his universal poem. The situation is resolved when the great writer decides to check the hoax for himself… and finds the prodigy, which he describes to us in his exuberant prose:

\begin{quotation}
Each thing (a mirror's face, let us say) was infinite things, since I distinctly saw it from every angle of the universe. I saw the teeming sea; I saw daybreak and nightfall; I saw the multitudes of America; […] I saw all the mirrors on earth and none of them reflected me; […] I saw tigers, pistons, bison, tides, and armies; I saw all the ants on the planet; […] I saw in the drawer of a writing table (and the handwriting made me tremble) unbelievable, obscene, detailed letters, which Beatriz had written to Carlos Argentino…
\end{quotation}

The enumeration of the marvels that the Aleph shows is the opposite pole to the botched sketches the dreadful cousin has written and, alas, the object of the erotic ravings of a Beatriz whom the point that contains all points strips, with one brutal swipe of the hand, of her place in Paradise, turning her instead into a woman of flesh and blood, whose absence, because of that detailed (and obscene) carnality, hurts even more. The parody is over, and the journey through the looking glass is tinged with pain.

And yet, the prodigy that Borges names with the first letter of the Hebrew alphabet ---the letter that contains all the others--- admits a reading that the writer himself, a citizen of his time, could not fully formulate but perhaps obscurely grasped: a point that contains all points describes exactly the infinitesimal instant that precedes the Big Bang, the great explosion that created the universe some 13.8 billion years ago. From there, from that point, that zero, that place of the orb that contains all the places of the orb, emanate matter and antimatter (the latter disappears at once, because of a traitor whose trail the writer of these lines has followed throughout his whole career), leptons and quarks that are confined in protons that form nuclei of hydrogen and then of helium and then of lithium and later beryllium… and then stars that explode in supernovae, which create solar systems where (inexplicably) life appears ---or at least that is what happens in one of them---: bacteria that, almost by magic, beget eukaryotes, eukaryotes that beget algae, algae that beget nautiluses, and centipedes, jellyfish and sharks, ferns and carnivorous plants, tyrannosaurs and crocodiles, sabre-toothed tigers, horses, unicorns and butterflies. And at some moment, blindly, there also appears a kind of mad monkey, capable of inventing literature to silence the noise inside its head, poetry to soothe its fear of death, science to try to understand the world.

But all of that happens afterwards. The Big Bang creates the universe, but the Aleph contains it in the infinitesimal time that precedes it. And that infinitesimal instant, physicists say, is a fluctuation of the quantum vacuum.

Two tremendous words. Vacuum. Quantum.

In the classical vacuum there is nothing. The theory of relativity tells us that matter weaves the very structure of space-time, so the vacuum, which forbids matter, cannot contain space or time either. It is an absolute absence.

But the laws of quantum mechanics assure us that such a vacuum does not exist. They preach that in the fabric of the cosmos, even in the absence of all matter, pairs of particles and antiparticles form and are destroyed, incessantly, in an invisible dance, hidden from the observer by the implacable uncertainty principle. And that phenomenon can fluctuate, generate a bubble, provoke a Big Bang, perhaps with a probability so small that it would be impossible for it to happen, unless one has all of eternity for it.

And so, eventually, perhaps in an inconceivable Multiverse, bubbles appear, initiated by that rebelliousness of the quantum vacuum, which form an Aleph, which in turn begets the orb.

Towards the end of his masterful story, Borges reckons that the Aleph that existed in the cellar of the poetaster Carlos Argentino was, possibly, a false Aleph. It is easy to believe that he is right and to venture that the authentic Aleph disappeared billions of years ago, giving rise to all things.

Or perhaps it is not so. Perhaps the quantum vacuum keeps fluctuating and other universes are continually forming, breaking away from ours like soap bubbles blown through the hollow pen of a mischievous God. Perhaps we have all seen an Aleph without knowing it, because that vision, too overwhelming, has been erased from our minds, as the features of Beatriz Viterbo have already been erased, forever.

\end{document}
